\section{Even and odd channels as diagnostic residuals}
We want to remove the legitimate symmetric magnetic structure while retaining
components that the model forbids. Compare points at equal $i_d$, temperature,
and signed speed, with opposite equal-magnitude $i_q$. Define
\begin{equation}
 \evenq[f](i_d,i_q)=\frac{f(i_d,i_q)+f(i_d,-i_q)}{2},\qquad
 \oddq[f](i_d,i_q)=\frac{f(i_d,i_q)-f(i_d,-i_q)}{2}.
 \label{eq:parity-operators}
\end{equation}
Addition retains the unchanged part; subtraction retains the reversing part.
These are complementary projections: $f=\evenq[f]+\oddq[f]$,
$\evenq^2=\evenq$, $\oddq^2=\oddq$, and $\evenq\oddq=0$. Pairing therefore discards the
allowed component in each diagnostic channel, not saturation as such.
\begin{figure}[tb]
\centering
\begin{tikzpicture}
\begin{axis}[diagnostic axis,title={Reflect the same curve},ylabel={$f(y)$},
 xmin=-1,xmax=1,ymin=.5,ymax=1.6]
 \addplot[reference quantity,strong line,domain=-1:1]{1+.2*x^2+.3*x};
 \addlegendentry{$f(y)$}
 \addplot[estimated quantity,alternative pattern,standard line,domain=-1:1]{1+.2*x^2-.3*x};
 \addlegendentry{$f(-y)$}
\end{axis}
\begin{axis}[diagnostic axis,at={(.5\textwidth,0)},anchor=south west,
 title={Average and half-difference},ylabel={Component},xmin=-1,xmax=1,ymin=-.4,ymax=1.4]
 \addplot[derived quantity,strong line,domain=-1:1]{1+.2*x^2};
 \addlegendentry{$\evenq[f]$}
 \addplot[torque quantity,standard line,domain=-1:1]{.3*x};
 \addlegendentry{$\oddq[f]$}
\end{axis}
\end{tikzpicture}
\caption{Reflection separates an illustrative curve
$f(y)=1+0.2y^2+0.3y$ into its allowed even part and its odd tilt.
The same projections applied to d and q flux select opposite diagnostic
channels. The schematic curve is not a second machine model.}
\label{fig:decomposition}
\end{figure}


For a correctly aligned ideal map, $\oddq[\psid]=0$ and $\evenq[\psiq]=0$.
Define measured residuals at the observed numerical labels by
\begin{equation}
 \boxed{r_d=\oddq[\hat\psid],\qquad r_q=\evenq[\hat\psiq].}
 \label{eq:residuals}
\end{equation}
For the independent positive-q representative $a>0$, this means explicitly
\begin{equation}
 r_d=\frac{\hat\psi_d(i_d,+a)-\hat\psi_d(i_d,-a)}{2},\qquad
 r_q=\frac{\hat\psi_q(i_d,+a)+\hat\psi_q(i_d,-a)}{2}.
\end{equation}
Using \cref{eq:R-error,eq:angle-sensitivity} gives
\begin{equation}
 \begin{bmatrix}r_d\\r_q\end{bmatrix}\approx
 \vect s_R\resistanceerror+\vect s_\gamma\angleerror+\vect r_{\rm other},
 \label{eq:fingerprint}
\end{equation}
where the ideal map and derivatives at the numerical argument define
\begin{equation}
 \vect s_R=\begin{bmatrix}-i_q/\omegae\\i_d/\omegae\end{bmatrix},\qquad
 \vect s_\gamma=\begin{bmatrix}
 -L_{dd}i_q+L_{dq}i_d+\psiq\\
 -L_{qd}i_q+L_{qq}i_d-\psid\end{bmatrix}.
 \label{eq:columns}
\end{equation}
Each term of $s_{\gamma,d}$ is odd: even times odd, odd times fixed $i_d$,
and odd flux. Each term of $s_{\gamma,q}$ is even: odd times odd, even times
fixed $i_d$, and even flux. Thus both errors enter precisely the forbidden
channels, but with different functional fingerprints.

For additive errors after coherent frame rotation, the other contribution
can be computed by applying $(\oddq,\evenq)$ componentwise to
\begin{equation}
 -\frac{\rotationgenerator\vect\varepsilon_u}{\omegae}
 +\left(\frac{\Rs}{\omegae}\rotationgenerator-\Ldiff\right)\vect\varepsilon_i
 -\frac{\flux}{\omegae}\varepsilon_\omega.\label{eq:other-map}
\end{equation}
For instance, a scalar speed error equal on both members of a pair multiplies
both ideal flux components by the same factor and preserves parity. Unequal
pair errors need not. A constant dq d-voltage offset creates a q-even term
$-\varepsilon_{u,d}/\omegae$, whereas a constant q-voltage offset shifts
d-even flux and is invisible to $r_d$. The coordinate representation of a
physical sensor offset must be determined before assigning these parities.

\paragraph{What pairing hides.}
An allowed-parity flux error can be large and yet yield zero residuals.
A nonzero residual indicates incompatibility with the specified symmetry,
not unique proof of resistance or angle error. Interpolation mismatch,
temperature mismatch, or real machine asymmetry can enter the same channels.
Use one independent representative per pair, conventionally $i_q>0$; keeping
both representatives doubles the same information and their correlated noise.

\subsection{A resistance-equivalent residual and its hypothesis}
To ask whether one constant resistance mismatch can explain the d channel,
remove its known current and speed factors. For ideal matched pairs,
\begin{equation}
 r_d=-\frac{\resistanceerror}{\omegae}a
 \quad\Longrightarrow\quad
 \boxed{\resistanceequivalent=-\omegae\frac{r_d}{a}.}
 \label{eq:resistance-equivalent}
\end{equation}
We call this the \emph{resistance-equivalent symmetry residual}. The notation
does not assert that the observed difference is a physical resistance error.
The null model $H_0$ is: a single constant used-minus-true resistance mismatch
is the sole source of the observed symmetry violation, with a correctly
aligned symmetric magnetic law and adequately matched measurements.
If that model holds, $\resistanceequivalent$ would be approximately constant
over operating point and speed. Its q-channel prediction is additionally
$r_q=\resistanceerror i_d/\omegae$; a flat d equivalent alone is insufficient.

In measured data the denominator is
\begin{equation}
 a_q=\frac{|i_{q,+}|+|i_{q,-}|}{2},\qquad
 \resistanceequivalent=-\omegae\frac{r_d}{a_q}.
 \label{eq:measured-denominator}
\end{equation}
For opposite-sign measured currents, the resistance contribution to the
half-difference is exactly $-\resistanceerror a_q/\omegae$.
This choice accounts for their differing magnitudes, but cannot cancel the
magnetic contribution caused by imperfectly mirrored measured current labels.
That distinction matters whenever measured pairs are matched or interpolated.

The equivalent transformation reveals departure from the constant-resistance
model, while hiding the original flux units and amplifying errors near zero q
current. A perturbation of $r_d$ contributes
$\delta\Delta R_{\mathrm{eq}}=-\omegae\delta r_d/a_q$ at fixed denominator.
The original $r_d$ is therefore retained and compared directly across speeds.
